% GENERATED FILE. Edit canonical lessons, then run npm run generate:books.
% canonical-source-hash: fnv1a64:2e936f9bef35f6f8
% canonical-lessons: GE-C333-widersprechen, GE-C333-widerspruch, GE-C333-einwand, GE-C333-einwenden, GE-C333-bestreiten

\chapter{Widersprechen}
\label{ch:ge-widersprechen}
\hlchaptermodality{333}

\begin{chapteropening}
\textbf{By the end of this chapter:} \emph{I can contradict someone, raise an objection and deny a claim.}

Complete the last lesson of chapter 333: I can contradict someone, raise an objection and deny a claim.
\end{chapteropening}

\section[widersprechen]{widersprechen — to contradict, to disagree with}

\label{lesson:GE-C333-widersprechen}



\begin{quote}
Before the new one: say the German for especially since, then the German for to explain in detail.
\end{quote}

\subsection*{You'll want to know: widersprechen}
\textbf{widersprechen} — \textquotedblleft{}to contradict, to disagree with\textquotedblright{}.

\begin{quote}
Ich muss dir widersprechen: Der Film war gar nicht langweilig.
\end{quote}

\begin{quote}
\emph{I have to contradict you: the film was not boring at all.}
\end{quote}

Inseparable, with the dative: \emph{ich widerspreche dir}, \emph{ich habe ihm widersprochen}. Conjugated like \emph{sprechen}: \emph{er widerspricht}.

\begin{cousinweb}
\emph{widersprechen} is inherited through Middle High German \emph{widersprechen} from Old High German \emph{widarsprehhan} (EtymDB), built from \emph{wider} \textquotedblleft{}against\textquotedblright{} and \emph{sprechen} \textquotedblleft{}to speak\textquotedblright{}: to speak against.
\end{cousinweb}

\begin{grammarlens}[title={What you've built}]
One more word for telling, explaining and describing at B1.
\end{grammarlens}

\subsection*{Your turn}
\begin{itemize}
\raggedright
  \item \emph{Say it:} \emph{widersprechen}
  \item \emph{Say it:} \emph{widersprechen}, once more
  \item \emph{Say it:} a sentence of your own with \emph{widersprechen}
  \item \emph{Recall:} say the German for especially since, then the German for to explain in detail, then say \emph{widersprechen} again
\end{itemize}

\begin{tcolorbox}[breakable,colback=teal!4,colframe=teal!35!black,title={Before you move on}]
What does widersprechen mean? (\textquotedblleft{}To contradict, to disagree with\textquotedblright{}.) Say it once more.
\end{tcolorbox}
\section[der Widerspruch]{der Widerspruch — the contradiction; the objection}

\label{lesson:GE-C333-widerspruch}



\begin{quote}
Before the new one: say the German for to explain in detail, then the German for to contradict.
\end{quote}

\subsection*{You'll want to know: der Widerspruch}
\textbf{der Widerspruch} — \textquotedblleft{}the contradiction; the objection\textquotedblright{}.

\begin{quote}
In seiner Geschichte gibt es einen klaren Widerspruch.
\end{quote}

\begin{quote}
\emph{There is a clear contradiction in his story.}
\end{quote}

Plural: \emph{die Widersprüche}. \emph{Ohne Widerspruch} — without objection.

\begin{cousinweb}
\emph{Widerspruch} is the noun that goes with \emph{widersprechen}, built from \emph{wider} \textquotedblleft{}against\textquotedblright{} and \emph{Spruch} \textquotedblleft{}saying\textquotedblright{}, from the family of \emph{sprechen} (EtymDB).
\end{cousinweb}

\begin{grammarlens}[title={What you've built}]
One more word for telling, explaining and describing at B1.
\end{grammarlens}

\subsection*{Your turn}
\begin{itemize}
\raggedright
  \item \emph{Say it:} \emph{der Widerspruch}
  \item \emph{Say it:} \emph{der Widerspruch}, once more
  \item \emph{Say it:} a sentence of your own with \emph{der Widerspruch}
  \item \emph{Recall:} say the German for to explain in detail, then the German for to contradict, then say \emph{der Widerspruch} again
\end{itemize}

\begin{tcolorbox}[breakable,colback=teal!4,colframe=teal!35!black,title={Before you move on}]
What does der Widerspruch mean? (\textquotedblleft{}The contradiction; the objection\textquotedblright{}.) Say it once more.
\end{tcolorbox}
\section[der Einwand]{der Einwand — the objection}

\label{lesson:GE-C333-einwand}



\begin{quote}
Before the new one: say the German for to contradict, then the German for the contradiction.
\end{quote}

\subsection*{You'll want to know: der Einwand}
\textbf{der Einwand} — \textquotedblleft{}the objection\textquotedblright{}.

\begin{quote}
Mein Vater hatte einen Einwand: Die Reise war ihm zu teuer.
\end{quote}

\begin{quote}
\emph{My father had an objection: the trip was too expensive for him.}
\end{quote}

Plural: \emph{die Einwände}. \emph{Einen Einwand gegen etwas haben} — to have an objection to something.

\begin{cousinweb}
\emph{Einwand} is formed from \emph{einwenden} \textquotedblleft{}to object\textquotedblright{} (EtymDB), built on \emph{wenden} \textquotedblleft{}to turn\textquotedblright{}, inherited from Old High German \emph{wenten} and Proto-Germanic \emph{wandijaną} (EtymDB): something you turn against an idea.
\end{cousinweb}

\begin{grammarlens}[title={What you've built}]
One more word for telling, explaining and describing at B1.
\end{grammarlens}

\subsection*{Your turn}
\begin{itemize}
\raggedright
  \item \emph{Say it:} \emph{der Einwand}
  \item \emph{Say it:} \emph{der Einwand}, once more
  \item \emph{Say it:} a sentence of your own with \emph{der Einwand}
  \item \emph{Recall:} say the German for to contradict, then the German for the contradiction, then say \emph{der Einwand} again
\end{itemize}

\begin{tcolorbox}[breakable,colback=teal!4,colframe=teal!35!black,title={Before you move on}]
What does der Einwand mean? (\textquotedblleft{}The objection\textquotedblright{}.) Say it once more.
\end{tcolorbox}
\section[einwenden]{einwenden — to object, to say against}

\label{lesson:GE-C333-einwenden}



\begin{quote}
Before the new one: say the German for the contradiction, then the German for the objection.
\end{quote}

\subsection*{You'll want to know: einwenden}
\textbf{einwenden} — \textquotedblleft{}to object, to say against\textquotedblright{}.

\begin{quote}
Sie wendet ein, dass der Weg zu weit ist.
\end{quote}

\begin{quote}
\emph{She objects that the way is too far.}
\end{quote}

Separable: \emph{sie wendet ein}. Perfekt: \emph{sie hat eingewandt} or \emph{eingewendet}; both are correct.

\begin{cousinweb}
\emph{einwenden} is built from \emph{ein} and \emph{wenden} \textquotedblleft{}to turn\textquotedblright{}, inherited from Old High German \emph{wenten} and Proto-Germanic \emph{wandijaną} (EtymDB).
\end{cousinweb}

\begin{grammarlens}[title={What you've built}]
One more word for telling, explaining and describing at B1.
\end{grammarlens}

\subsection*{Your turn}
\begin{itemize}
\raggedright
  \item \emph{Say it:} \emph{einwenden}
  \item \emph{Say it:} \emph{einwenden}, once more
  \item \emph{Say it:} a sentence of your own with \emph{einwenden}
  \item \emph{Recall:} say the German for the contradiction, then the German for the objection, then say \emph{einwenden} again
\end{itemize}

\begin{tcolorbox}[breakable,colback=teal!4,colframe=teal!35!black,title={Before you move on}]
What does einwenden mean? (\textquotedblleft{}To object, to say against\textquotedblright{}.) Say it once more.
\end{tcolorbox}
\section[bestreiten]{bestreiten — to deny, to dispute}

\label{lesson:GE-C333-bestreiten}



\begin{quote}
Before the new one: say the German for the objection, then the German for to object.
\end{quote}

\subsection*{You'll want to know: bestreiten}
\textbf{bestreiten} — \textquotedblleft{}to deny, to dispute\textquotedblright{}.

\begin{quote}
Der Junge bestreitet, dass er das Fenster kaputt gemacht hat.
\end{quote}

\begin{quote}
\emph{The boy denies that he broke the window.}
\end{quote}

Inseparable and strong: \emph{er bestreitet, er hat bestritten}. \emph{Das kann man nicht bestreiten} — that cannot be denied.

\begin{cousinweb}
\emph{bestreiten} is inherited from Middle High German \emph{bestrīten} \textquotedblleft{}to fight against\textquotedblright{} (EtymDB), built on \emph{strīten}, today \emph{streiten} \textquotedblleft{}to quarrel\textquotedblright{}: you fight a claim.
\end{cousinweb}

\begin{grammarlens}[title={What you've built}]
One more word for telling, explaining and describing at B1.
\end{grammarlens}

\subsection*{Your turn}
\begin{itemize}
\raggedright
  \item \emph{Say it:} \emph{bestreiten}
  \item \emph{Say it:} \emph{bestreiten}, once more
  \item \emph{Say it:} a sentence of your own with \emph{bestreiten}
  \item \emph{Recall:} say the German for the objection, then the German for to object, then say \emph{bestreiten} again
\end{itemize}

\begin{tcolorbox}[breakable,colback=teal!4,colframe=teal!35!black,title={Before you move on}]
What does bestreiten mean? (\textquotedblleft{}To deny, to dispute\textquotedblright{}.) Say it once more.
\end{tcolorbox}
